\documentclass[../principal.tex]{subfiles}
\graphicspath{{\subfix{../imagenes/}}}

\begin{document}

\chapter{Clases}

Las clases son estructuras de código que son usadas en el paradigma de programación llamado programación orientada a objetos.

Las clases nos permiten hacer moldes reutilizables, como si creáramos nuestra propio tipo de variable.

Por ejemplo si pensamos en las variables enteras, todas ellas tienen comportamientos similares que las hacen variables enteras, pero tambiéncada instancia tiene su identidad, al tener su propio nombre y su propio valor.

Otro comportamiento importante es que podemos operar y relacionar números enteros con otros números enteros mediante álgebra, o mediante operadores lógicos.

\section{Declaración de una clase}

La declaración de una clase en C++ sigue la siguiente estructura.

\begin{lstlisting}[language=C++]
class NombreDeLaClase {
    public:
      // atributos

      // constructores

      // métodos
};
\end{lstlisting}

Se usa la palabra reservada \lstinline{class}, luego se escribe el nombre de la clase, que sigue la convención de tener un nombre con notación Pascal.

Notar que por ahora estamos declarando todo como \lstinline{public} dentro de la clase, pero lo revisitaremos más adelante.

Luego sus componentes internos están encerrados por llaves \{\} y finaliza con un punto y coma.

\section{Constructores}

\section{Atributos}

Los atributos son variables internas

\section{Métodos}


\newpage

\end{document}

